\chapter{Inscrição profissional de pessoa física}
 
\section*{O que é?}

É o registro necessário para a atuação profissional em Psicologia. O exercício profissional no estado de São Paulo estará autorizado após a geração do número de registro profissional.

\section*{Quem deve se inscrever?}

Profissionais graduadas/os em Psicologia, que tenham atuação no estado de São Paulo.

\section*{Como solicitar o serviço?}

Por meio do site do CRP SP, na guia Atendimento \> Pessoa Física.

\section*{Que informações ou documentos são necessários?}

\begin{itemize}
\item	 Documento de identificação (RG, CNH etc.);
\item	 Cadastro de Pessoa Física (CPF);
\item	 certidão de quitação eleitoral;
\item	 comprovante ou declaração de endereço;
\item	 diploma de formação em Psicologia (frente e verso) ou certidão de colação de grau;
\item	 certidão de casamento ou de união estável;
\item	 comprovante de quitação com o serviço militar (para homens até 45 anos).
\end{itemize}

\section*{Quais são as etapas para a realização desse serviço?}

\begin{enumerate}
\item	A/O psicóloga/o realiza o requerimento de inscrição on-line, pelo site do CRP SP;
\item	o CRP SP faz a análise da documentação;
\item	os boletos para pagamento da taxa de inscrição e da anuidade são emitidos;
\item	após o pagamento das taxas, o registro é deferido e uma declaração para atuação profissional é disponibilizada.
\end{enumerate}

\section*{Quanto custa?}

Os valores para 2024 ficaram estabelecidos assim:

\begin{itemize}
\item	pagamento à vista: R\$ 575,60;
\item	pagamento parcelado: 6 x R\$ 95,93.
\end{itemize}

\section*{Quanto tempo leva?}

Até 30 dias. 

\section*{Serviços correlatos}

\begin{itemize}
\item	Entrega carteira de identidade profissional (CIP) (Resolução CFP n. 02/2021);
\item	alteração de inscrição provisória para definitiva (Resolução CFP n. 03/2007);
\item	inscrição secundária (Resolução CFP n. 03/2007)
\item	transferência do registro profissional (Resolução CFP n. 03/2007);
\item	cancelamento da inscrição (Resolução CFP n. 03/2007);
\item	prorrogação do prazo de entrega do diploma  (Resolução CFP n. 03/2007);
\item	solicitação de isenção de anuidade (Resolução CFP n. 08/2023).
\end{itemize}

\section*{Como ter acesso ao atual status do serviço?}

Ao se fazer a solicitação é gerada uma senha de acesso. Com essa senha, é possível acompanhar o status da solicitação \textit{on-line}.

\section{Legislação relacionada}

\begin{itemize}
\item	Resolução CFP n. 03/2007 (Consolidação das Resoluções do Conselho Federal de Psicologia);
\item	Resolução CFP n. 20/2018, que determina a revisão e ampliação do Manual de procedimentos administrativos e financeiros (versão atualizada).
\end{itemize}

